\section{Methods and Materials}
\label{sec:methods}

\subsection{Data, Temporal Split and Supervised Components}
\label{sec:data}

We reuse the Workshop~1 (W1) sample of MovieLens 32M. Fifteen thousand users are drawn uniformly at random without replacement (seed 42), all their ratings are kept, and an iterative 10-core filter removes users and movies with fewer than ten ratings until both conditions hold. The sample keeps 2{,}368{,}980 ratings given by 14{,}998 users to 11{,}589 movies. The ratings are sorted by timestamp and split globally into the first 80\%, the next 10\% and the last 10\%, the chronological split of W1 and Workshop~2 (W2): the training split holds 1{,}895{,}184 ratings up to 18 August 2018, and the validation and test splits hold 236{,}898 ratings each, up to September 2020 and October 2023. A rating outside the training split is \emph{clean} when its user and its movie both appear in the training split. A FunkSVD model with $k = 20$ factors, learning rate 0.01, regularization 0.02 and 15 epochs of mini-batch stochastic gradient descent is trained on the training split only; its RMSE is 0.6954 on the training split, 0.7596 on the clean validation ratings and 0.8128 on the clean test ratings. The global mean $\mu$, the biases $b_u$ and $b_i$, the factors $\mathbf{p}_u$ and $\mathbf{q}_i$, the item popularity and the self-information novelty $-\log_2(n_i/N)$, where $n_i$ is the training count of movie $i$ and $N$ the size of the training split, all come from the training split, so every feature the agent observes predates every RL session.

\subsection{Sessions and Candidate Slates}
\label{sec:sessions}

An RL session contains the clean ratings of one user inside one period, in chronological order. A user is eligible when the period holds at least 20 clean ratings and the first 50 of them contain at least one like, where a like (a \emph{hit} when recommended) is a real rating of at least 4.0. The candidate slate of a session holds those first 50 ratings: it is chosen by timestamp alone and then ordered by the static FunkSVD score, so the slot order never depends on the ratings, and shorter sessions are padded with masked slots. The validation period yields 284 eligible users, split at random into 227 RL training users and 57 RL validation users. The test period yields 177 RL test users, the same held-out future on which W1 and W2 report; 98 of them also have an earlier RL training session, which mirrors the temporal protocol of the supervised workshops. Training and test slates hold 42.81 and 43.77 real slots on average, of which 20.93 and 19.89 are likes. The environment uses the ratings only to emit rewards and to compute metrics; no rating or label enters the observation.

\subsection{MDP Formulation}
\label{sec:mdp}

An episode replays the session of one user $u$. At step $t$ the agent picks a slot $a_t \in \mathcal{A} = \{1,\dots,50\}$ and observes the real rating $y_{a_t}$ of the movie in that slot.

\emph{State.} The observation $s_t \in [-5,5]^{250}$ concatenates four 50-dimensional slot blocks, two 20-dimensional vectors and ten scalars. For slot $a$, the blocks hold the cosine $c_t(a) = \hat{\mathbf{q}}_a^{\top}\bar{\mathbf{s}}_t$ between the unit FunkSVD factor of its movie and the session vector $\bar{\mathbf{s}}_t$, the FunkSVD score $\mu + b_u + b_i + \mathbf{p}_u^{\top}\mathbf{q}_i$ min--max normalized within the slate, the novelty min--max normalized over the catalog, $\tilde{\nu}_a \in [0,1]$, and an unavailable flag $m_t(a)$ that equals one for recommended and padding slots. The two vectors are the normalized user factor $\hat{\mathbf{p}}_u$ and $\bar{\mathbf{s}}_t$. The scalars are $t/H$, $(H-t)/H$, the hit rate so far, the last reward, the number of consecutive misses divided by five, the mean novelty and the genre coverage of the movies recommended so far, the user bias, and the maximum cosine and mean FunkSVD score of the available slots.

\emph{Transitions.} The session vector starts at $\bar{\mathbf{s}}_0 = \hat{\mathbf{p}}_u$ and follows a Rocchio-style relevance-feedback update that moves it toward liked movies and away from missed ones:
\begin{equation}
\bar{\mathbf{s}}_{t+1} = \frac{\mathbf{v}_t}{\lVert \mathbf{v}_t \rVert},\quad
\mathbf{v}_t =
\begin{cases}
0.50\,\bar{\mathbf{s}}_t + 0.30\,\hat{\mathbf{q}}_{a_t} & \text{if } y_{a_t} \ge 4,\\
0.75\,\bar{\mathbf{s}}_t - 0.20\,\hat{\mathbf{q}}_{a_t} & \text{if } y_{a_t} < 4.
\end{cases}
\label{eq:session}
\end{equation}
The chosen slot becomes unavailable, and a repeated or padding slot leaves $\bar{\mathbf{s}}_t$ unchanged.

\emph{Reward.} With $g_t$ the fraction of the movie's genres not yet covered in the episode,
\begin{equation}
r_t =
\begin{cases}
-1.5 & \text{if } m_t(a_t) = 1,\\
1 + 0.25\,(y_{a_t} - 4) + 0.20\,\tilde{\nu}_{a_t} + 0.10\,g_t & \text{if } y_{a_t} \ge 4,\\
-0.4 & \text{otherwise}.
\end{cases}
\label{eq:reward}
\end{equation}

\emph{Termination.} The user abandons the session, a terminal state, after five consecutive misses, repeats included. Otherwise the episode is truncated at the horizon: $H = 10$ for training and validation, and $K \in \{10, 20\}$ at test time.

\emph{Justification and assumptions.} The formulation keeps the three ingredients that make the task sequential. Feedback changes the observation through \eqref{eq:session}, so a miss lowers the cosine of similar movies; abandonment makes a streak of misses cost more than its immediate penalties, which gives planning a role; and the unavailable flags rule out recommending the same movie again. The novelty and genre bonuses are small and paid only on hits, so they cannot be earned with disliked long-tail movies, and every reward lies between $-1.5$ and about $1.5$. The simulator rests on two strong assumptions. The ratings are fixed, so the order of the recommendations changes the return only through abandonment and the genre bonus: the session vector changes what the agent knows, not what the user wants. The agent can also recommend only movies the user rated in that period, a restriction inherent to learning from logged data~\cite{levine2020offline} that hides how the user would react to any other movie. Both assumptions bound how much a sequential policy can gain over a good static ranking.

\subsection{Environment Validation}
\label{sec:envval}

The environment, \texttt{MovieLensRecEnv}, follows the Gymnasium API~\cite{towers2024gymnasium} with a \texttt{Discrete(50)} action space, a \texttt{Box(250)} observation space and the five-value \texttt{step} output, and returns the valid-action mask in its \texttt{info} dictionary. Before any training, a uniform random policy over all 50 slots, without masking, runs 100 episodes on randomly drawn test sessions with $H = 10$. The run checks that every observation lies inside the observation space and every reward is finite, that both abandonment and truncation occur, and that every action index is sampled; it also measures the return, hit rate and repeat rate that any learned policy must exceed.

\subsection{Agent and Q-Network Architectures}
\label{sec:agent}

The agent is a Dueling Double DQN~\cite{mnih2015human,van2016deep,wang2016dueling} implemented in PyTorch. Transitions $(s_t, a_t, r_t, s_{t+1}, d_t, \mathbf{m}_{t+1})$ enter a replay buffer of 50{,}000 transitions, where $d_t$ marks the end of the episode and $\mathbf{m}_{t+1}$ the unavailable slots of the next state. Because $t/H$ is part of the state, the horizon belongs to the MDP and truncation is treated as terminal. After every environment step, once the buffer holds 128 transitions, one mini-batch of 128 updates the online network $Q_\theta$ with the Huber loss toward the Double DQN target
\begin{equation}
\begin{aligned}
z_t &= r_t + \gamma\,(1-d_t)\,Q_{\bar{\theta}}(s_{t+1}, a^{*}),\\
a^{*} &= \mathop{\arg\max}\nolimits_{a:\,m_{t+1}(a)=0}\, Q_{\theta}(s_{t+1}, a),
\end{aligned}
\label{eq:target}
\end{equation}
in which the online network chooses the next action among the available slots and the target network $Q_{\bar{\theta}}$ evaluates it. The optimizer is AdamW with learning rate $\eta$ and decoupled weight decay $\lambda$ (the L2 term of the search), gradient norms are clipped at 1.0, and the target network follows the Polyak average $\bar{\theta} \leftarrow 0.95\,\bar{\theta} + 0.05\,\theta$ after every update. Exploration is $\epsilon$-greedy over the available slots: $\epsilon$ starts at 1.0 and is multiplied by $\delta_\epsilon$ after every episode down to 0.05. Validation and test episodes use the greedy policy.

Both Q-networks are dueling, $Q(s,a) = V(s) + A(s,a) - \frac{1}{50}\sum_{a'} A(s,a')$. The \emph{flat} network is an MLP over the whole state with two hidden layers of 256 units (LayerNorm and ReLU) followed by value and advantage streams of 128 units; each slot has its own advantage output (203{,}443 parameters). The \emph{slot-shared} network scores every slot with one MLP. A slot is described by five features: its cosine, FunkSVD score, novelty, unavailable flag and rank in the slate. A context MLP embeds the ten scalars into 64 units, the shared scorer maps the slot features and the context through layers of 128 and 64 units to the advantage of the slot, and the value head reads the context and the mean slot embedding (26{,}434 parameters). This network deliberately ignores the raw user factor and session vector: with 227 training users, those 40 inputs identify each user and let a network memorize training sessions instead of learning a scoring rule. Personalization still enters through the FunkSVD score and the session cosine of every slot, and sharing the scorer lets one transition inform the value of every slot with similar features.

\subsection{Hyperparameter Search}
\label{sec:search}

The search varies one hyperparameter at a time around a base configuration with $\eta = 5\times10^{-4}$, $\gamma = 0.95$, $\delta_\epsilon = 0.995$ and $\lambda = 10^{-5}$: $\eta \in \{10^{-3}, 10^{-4}\}$, $\gamma \in \{0.90, 0.99\}$, $\delta_\epsilon \in \{0.990, 0.998\}$ and $\lambda = 10^{-4}$. The eight configurations are crossed with the two architectures, which gives 16 runs of 800 episodes each with seed 42; each training episode draws one of the 227 training sessions at random. Before training and every 40 episodes, the greedy policy plays each of the 57 validation sessions once. We select the run with the highest smoothed validation return, the mean of its last five evaluations, which is less sensitive to one lucky evaluation than the maximum. The selected configuration is retrained from scratch for 1{,}500 episodes, and its checkpoint with the best validation return becomes the champion. A copy with $\gamma = 0$ and the same network, data, budget and checkpoint rule turns the agent into a neural contextual bandit that maximizes only the immediate reward, so the gap between the two measures the value of multi-step planning. The three exploration schedules reach $\epsilon = 0.05$ after about 300, 600 and 1{,}500 episodes, so the slow schedule still explores with $\epsilon \approx 0.20$ when the 800-episode budget ends. The test sessions are used once, after every choice has been made.

\subsection{Baselines}
\label{sec:baselines}

Every baseline acts greedily inside the same environment and never selects a recommended or padding slot. Random picks an available slot uniformly (seed 42). Popularity ranks slots by training popularity, and Bias by $\mu + b_u + b_i$. Greedy FunkSVD (W1) ranks by the static FunkSVD score, which is the slate order. Greedy HybridNeuMF (W2) ranks by the rating that the official W2 checkpoint predicts for every slot. Before scoring, the W2 pipeline is rebuilt and must reproduce the official W2 test RMSE, MAE, ROC-AUC and ranking metrics. All 461 RL users and all 19{,}811 real slots map to W2 indices, no slot involves a movie unseen in W2 training (0 cold slots), no slot's rating belongs to the W2 training data (0 slots), and every slot matches a W2 validation or test rating, so the W2 model scores only ratings it never saw. The Myopic DQN is the $\gamma = 0$ ablation.

\subsection{Evaluation Protocol}
\label{sec:protocol}

Every policy plays one episode on each of the 177 test sessions at $K = 10$ and at $K = 20$, so all policies face the same users. Precision@$K$ (P) counts hits among the first $K$ steps, and steps lost to abandonment count as misses. Recall, NDCG and MAP treat the session's liked slots as the relevant set, with binary gains and an ideal list of $\min(K, \text{likes})$ hits; HR@$K$ is the fraction of sessions with at least one hit. Coverage divides the number of distinct recommended movies by the number of distinct movies in the evaluated slates, and novelty is the mean self-information of the recommended movies in bits. For the main pairs of policies we report the mean per-user difference at $K = 10$ with a 95\% paired bootstrap interval (10{,}000 resamples) and a two-sided Wilcoxon signed-rank test. Sample efficiency is measured on the validation sessions: the champion's validation curve is plotted against cumulative environment steps and compared with the validation return and NDCG@10 of each baseline on the same 57 users. The W1 and W2 metrics from their own reports rank each test user's liked movies among sampled unrated movies, so we quote them as context and not as one-to-one comparisons.

\subsection{Software Tests}
\label{sec:tests}

A pytest suite checks the Gymnasium API, the mask, the penalties, abandonment and truncation. Leakage tests recompute every feature from the training-split FunkSVD parameters, verify that every session lies after the training cut, and require the initial session cosine to separate liked from disliked slots with an AUC below 0.75 (scikit-learn~\cite{pedregosa2011scikit}). Further tests check the alignment of the W2 scores, execute the notebook, and require every reported number to be traceable to a JSON result file.
